\let\negmedspace\undefined
\let\negthickspace\undefined
\documentclass[journal]{IEEEtran}
\usepackage[a5paper, margin=10mm, onecolumn]{geometry}
\usepackage{tfrupee}

\setlength{\headheight}{1cm}
\setlength{\headsep}{0mm}

\usepackage{gvv-book}
\usepackage{gvv}
\usepackage{cite}
\usepackage{amsmath,amssymb,amsfonts,amsthm}
\usepackage{algorithmic}
\usepackage{graphicx}
\usepackage{textcomp}
\usepackage{xcolor}
\usepackage{listings}
\usepackage{enumitem}
\usepackage{mathtools}
\usepackage{gensymb}
\usepackage{comment}
\usepackage[breaklinks=true]{hyperref}
\usepackage{tkz-euclide}
\usepackage{array}
\usepackage{circuitikz}

\tikzstyle{block} = [rectangle, draw, fill=blue!20, 
    text width=4em, text centered, rounded corners, minimum height=3em]
\tikzstyle{sum} = [draw, fill=blue!10, circle, minimum size=1cm, node distance=1.5cm]
\tikzstyle{input} = [coordinate]
\tikzstyle{output} = [coordinate]

\begin{document}

\bibliographystyle{IEEEtran}
\vspace{3cm}

\title{5.8.7}
\author{EE25BTECH11009 - Anshu kumar Ram}
\maketitle
{\let\newpage\relax\maketitle}

\renewcommand{\thefigure}{\theenumi}
\renewcommand{\thetable}{\theenumi}
\setlength{\intextsep}{10pt}

\numberwithin{equation}{enumi}
\numberwithin{figure}{enumi}
\renewcommand{\thetable}{\theenumi}

\textbf{Question}:\\
Yash scored 40 marks in a test, getting 3 marks for each right answer and losing 1
mark for each wrong answer. Had 4 marks been awarded for each correct answer
and 2 marks been deducted for each incorrect answer, then Yash would have scored
50 marks. How many questions were there in the test?

\solution \\

Let the number of right and wrong answers be $x$ and $y$ respectively.
\begin{align}
\myvec{3 & -1}\myvec{x \\ y} = 40
\end{align}
\begin{align}
\myvec{4 & -2}\myvec{x \\ y} = 50
\end{align}

These can be combined to give the matrix equation
\begin{align}
\myvec{3 & -1 \\ 4 & -2}\myvec{x \\ y} = \myvec{40 \\ 50}
\end{align}

This gives the augmented matrix
\begin{align}
\augvec{2}{1}{3 & -1 & 40 \\ 4 & -2 & 50} 
\xleftrightarrow{R_2 \leftarrow R_2 - \frac{4}{3}R_1} 
\augvec{2}{1}{3 & -1 & 40 \\ 0 & -\frac{2}{3} & -\frac{10}{3}}
\end{align}

This gives the following values of $x$ and $y$:
\begin{align}
\myvec{x \\ y} = \myvec{15 \\ 5}
\end{align}

Thus, the number of questions in the test is 
\begin{align}
x + y = 20
\end{align}


\begin{figure}[h!]
    \centering
    \includegraphics[height=0.5\textheight, keepaspectratio]{figs/Figure_1.png}
    \label{figure_1}
\end{figure}

\end{document}
